\section*{Capítulo \ref{ch:b1:plate-kinematics} --- Cinemática de placas sobre uma esfera}

\begin{solution}{exo:b1:plate-kinematics:1}
(a) $\qty{0.210}{\degree\per\Myr} = 0.210\times\pi/180 = \qty{3.665e-3}{rad}$
por milhão de anos, isto é \qty{3.665e-9}{rad.yr^{-1}}; com
$\qty{1}{yr} = \qty{3.156e7}{s}$, \qty{1.161e-16}{rad.s^{-1}}. (b) $360/0.210 =
\qty{1714}{\Myr}$, mais de um terço da idade da Terra. (c) $\omega R =
\num{3.665e-9}\times\qty{6.371e9}{mm} = \qty{23.35}{mm.yr^{-1}}$.
\end{solution}

\begin{solution}{exo:b1:plate-kinematics:2}
$v = \omega R\sin\Delta = 23.35\sin\Delta$: \qty{11.7}{mm.yr^{-1}} a \ang{30},
\qty{20.2}{mm.yr^{-1}} a \ang{60} e \qty{23.35}{mm.yr^{-1}} a \ang{90}. O
espalhamento é mais rápido onde a fronteira está mais longe do polo, até
\ang{90}: entre a Eurásia e a América do Norte, na extremidade sul da
fronteira, perto dos Açores, a \ang{78} do polo.
\end{solution}

\begin{solution}{exo:b1:plate-kinematics:3}
O subcron durou $1.072 - 0.988 = \qty{0.084}{\Myr}$, e uma taxa de metade em
\unit{mm.yr^{-1}} é também uma taxa em \unit{km\per\Myr}. (a) $50\times0.084 =
\qty{4.2}{km}$ a \qty{50}{mm.yr^{-1}}; $10\times0.084 = \qty{0.84}{km}$ a
\qty{10}{mm.yr^{-1}}. (b) A listra atravessa o eixo:
$2\times10\times0.781 = \qty{15.6}{km}$.
\end{solution}

\begin{solution}{exo:b1:plate-kinematics:4}
As taxas de metade são $66/2.581 = \qty{25.6}{mm.yr^{-1}}$ para oeste e
$64/2.581 = \qty{24.8}{mm.yr^{-1}}$ para leste; a taxa completa é
$130/2.581 = \qty{50.4}{mm.yr^{-1}}$. O
espalhamento é quase simétrico: o flanco oeste cresceu cerca de
\qty{3}{\percent} mais depressa.
\end{solution}

\begin{solution}{exo:b1:plate-kinematics:5}
$\cos\Delta = \sin\ang{61.8}\sin\ang{85} + \cos\ang{61.8}\cos\ang{85}
\cos(\ang{0} - \ang{139.6}) = 0.8813\times0.9962 + 0.4726\times0.08716\times
(-0.7615) = 0.87795 - 0.03137 = 0.84658$, logo $\Delta = \ang{32.16}$ e
$v = 23.35\sin\ang{32.16} = \qty{12.4}{mm.yr^{-1}}$. Isto está um pouco acima das
\qtyrange{11}{12}{mm.yr^{-1}} medidas, mas o \omterm{def:b1:plate-kinematics:euler-pole}{vetor de Euler} é um único melhor
ajuste a 453 taxas, das quais as taxas se afastam em \qty{0.7}{mm.yr^{-1}} em
média, e taxas individuais perto da lenta dorsal de Gakkel estão entre as menos
precisas.
\end{solution}

\begin{solution}{exo:b1:plate-kinematics:6}
(a) $\cos\Delta = \sin\ang{61.8}\sin\ang{39.44} + \cos\ang{61.8}\cos\ang{39.44}
\cos(\ang{169.43}) = 0.20111$: $\Delta = \ang{78.40}$. (b) $v = 23.35
\sin\ang{78.40} = \qty{22.87}{mm.yr^{-1}}$, a \qty{23}{mm.yr^{-1}} medida
logo ao norte da junção. (c) Com a fórmula do método, o numerador
é $\sin\ang{169.43}\cos\ang{61.8} = 0.08668$ e o denominador
$\cos\ang{39.44}\sin\ang{61.8} - \sin\ang{39.44}\cos\ang{61.8}\cos\ang{169.43}
= 0.97573$: o polo é visto em $\beta = \ang{5.08}$, quase exatamente ao norte, e
a Eurásia se move em direção a $\ang{95.1}$. (d) As transformantes são paralelas ao
movimento relativo: elas têm \omterm{def:b1:geological-maps:attitude}{direção} de cerca de \ang{095}, quase leste--oeste.
\end{solution}

\begin{solution}{exo:b1:plate-kinematics:7}
(a) Em componentes, $\vec\omega = \omega(\cos\varphi_p\cos\lambda_p, \cos\varphi_p\sin\lambda_p,
\sin\varphi_p)$: $\vec\omega_{\mathrm{EU\,NA}} = (-0.0756, 0.0643, 0.1851)$ e, de modo análogo,
$\vec\omega_{\mathrm{NB\,NA}} = (0.0333, 0.0282, 0.2289)$, em
\unit{\degree\per\Myr}. (b) $\vec\omega_{\mathrm{NB\,EU}} =
\vec\omega_{\mathrm{NB\,NA}} - \vec\omega_{\mathrm{EU\,NA}} = (0.1089, -0.0361,
0.0438)$, de comprimento
\[
  \sqrt{0.011859 + 0.001303 + 0.001918} = \qty{0.1228}{\degree\per\Myr}.
\]
Seu polo está na \omterm{def:b1:spherical-astronomy:geographic}{latitude}
$\arcsin(0.0438/0.1228) = \ang{20.9}$\,N e \omterm{def:b1:spherical-astronomy:geographic}{longitude}
$\arctan(-0.0361/0.1089) = \ang{-18.3}$, isto é \ang{18.3}\,W (a
componente $x$ sendo positiva, nenhum \ang{180} é somado). (c) A fronteira entre
a Núbia e a Eurásia oferece poucas medidas diretas: entre os Açores e
Gibraltar, apenas os \omterm{def:b1:spherical-astronomy:horizontal}{azimutes} da falha de Gloria. Os dois movimentos de espalhamento
através do Atlântico são medidos por centenas de taxas, e sua diferença
herda sua precisão.
\end{solution}

\begin{solution}{exo:b1:plate-kinematics:8}
(a) Em torno de $\vec e_z$ por \ang{90}: $\cos\ang{90} = 0$, $\vec e_z\cdot\vec e_x = 0$
e $\vec e_z\times\vec e_x = \vec e_y$, logo $\vec e_x$ vai para $\vec e_y$. Depois
em torno de $\vec e_x$ por \ang{90}: $\vec e_x\cdot\vec e_y = 0$ e
$\vec e_x\times\vec e_y = \vec e_z$, logo $\vec e_y$ vai para $\vec e_z$: o ponto
termina no polo norte. (b) Em torno de $\vec e_x$ primeiro: $\vec e_x$ está sobre o eixo,
$\vec e_x\times\vec e_x = \vec 0$ e $(1 - \cos\ang{90})(\vec e_x\cdot\vec e_x)
\vec e_x = \vec e_x$, logo não se move; depois, em torno de $\vec e_z$, vai para
$\vec e_y$, a \ang{0}\,N, \ang{90}\,L. (c) Os resultados diferem: \omterm{def:b1:plate-kinematics:finite-rotation}{rotações finitas}
não comutam. Reconstruções sucessivas devem ser compostas por
multiplicação das matrizes de rotação na ordem certa, nunca somando ângulos
ou polos.
\end{solution}

\begin{solution}{exo:b1:plate-kinematics:9}
(a) No espaço de velocidades ponha $a$ na origem, $b$ a \qty{30}{mm.yr^{-1}} para o
leste e $c$ a \qty{40}{mm.yr^{-1}} para o sul. (b) $\vec v_{CB} =
\vec v_{CA} - \vec v_{BA} = (0, -40) - (30, 0) = (-30, -40)$ (componentes leste, norte):
\qty{50}{mm.yr^{-1}} em direção ao \omterm{def:b1:spherical-astronomy:horizontal}{azimute}
$\arctan\bigl((-30)/(-40)\bigr) + \ang{180} = \ang{216.9}$, sudoeste. (c) A
dorsal é perpendicular a essa direção: ela tem \omterm{def:b1:geological-maps:attitude}{direção}
$\ang{216.9} - \ang{90} = \ang{126.9}$, isto é noroeste--sudeste.
\end{solution}

\begin{solution}{exo:b1:plate-kinematics:10}
(a) Na matriz de Rodrigues, a identidade contribui com $3\cos\alpha$, a matriz
$\vec n\,\vec n^{\mathsf T}$ com traço $n_x^2 + n_y^2 + n_z^2 = 1$, vezes
$1 - \cos\alpha$, e a matriz antissimétrica nada: $\operatorname{tr}R =
3\cos\alpha + 1 - \cos\alpha = 1 + 2\cos\alpha$. (b) Aqui $\operatorname{tr}R = 0$,
logo $\cos\alpha = -1/2$ e $\alpha = \ang{120}$. O eixo satisfaz $R\vec n =
\vec n$: $R(x, y, z) = (z, x, y)$, logo $x = y = z$ e $\vec n =
(1, 1, 1)/\sqrt3$, na \omterm{def:b1:spherical-astronomy:geographic}{latitude} $\arcsin(1/\sqrt3) = \ang{35.26}$\,N e
\omterm{def:b1:spherical-astronomy:geographic}{longitude} \ang{45}\,L. (c) $R\vec e_x = \vec e_y$, $R\vec e_y = \vec e_z$,
$R\vec e_z = \vec e_x$: os três eixos são permutados, como um terço de volta em torno
da diagonal. A fórmula de Rodrigues com $\alpha = +\ang{120}$ em torno de $\vec n$
de fato dá $\vec e_x\mapsto -\tfrac12\vec e_x + \tfrac12(1, 1, 1) +
\tfrac{\sqrt3}{2}\cdot\tfrac{1}{\sqrt3}(0, 1, -1) = (0, 1, 0)$, logo a rotação
é anti-horária vista de cima de \ang{35.26}\,N, \ang{45}\,L.
\end{solution}

\begin{solution}{exo:b1:plate-kinematics:11}
(a) Em uma transformante a velocidade relativa $\vec\omega\times\vec r$ está ao longo da
falha, logo paralela a $\vec t$; ela é perpendicular a $\vec\omega$, portanto
$\vec t\cdot\vec\omega = 0$. O plano perpendicular a $\vec t$ pelo
centro contém $\vec r$ (pois $\vec t$ é horizontal) e $\vec\omega$: ele
é o plano do \omterm{def:b1:spherical-astronomy:spherical-triangle}{círculo máximo} que passa pelo ponto, perpendicular à falha,
e o polo está sobre esse círculo. Dois desses planos se encontram ao longo da reta de
$\vec t_1\times\vec t_2$. (b) $\vec t_1 = (0.5436, 0.8367, -0.0673)$ na
junção e $\vec t_2 = (0.4942, 0.8639, -0.0968)$ na Islândia, logo
$\vec t_1\times\vec t_2 = (-0.02285, 0.01936, 0.05612)$, de comprimento $0.06362$.
Sua \omterm{def:b1:spherical-astronomy:geographic}{latitude} é $\arcsin(0.05612/0.06362) = \ang{61.9}$\,N e sua \omterm{def:b1:spherical-astronomy:geographic}{longitude}
$\ang{180} - \arctan(0.01936/0.02285) = \ang{139.7}$\,L (componente $x$ negativa):
o polo do capítulo. (c) Mal: os dois \omterm{def:b1:spherical-astronomy:spherical-triangle}{círculos máximos} se cruzam em um ângulo pequeno,
pois as direções diferem em apenas \ang{8} ao longo de \ang{26} da fronteira.
Mudar qualquer uma das direções em um grau move o polo em cerca de \ang{15}. Determinações
reais combinam muitas transformantes com taxas de espalhamento, cuja variação
ao longo da fronteira fixa a distância angular ao polo.
\end{solution}

\begin{solution}{exo:b1:plate-kinematics:12}
(a) $(94 - 4)/3 = \qty{30}{mm.yr^{-1}}$. (b) As diferenças separadas por um ano são
$24 - 4 = 20$, $25 - 5 = 20$, $26 - 6 = 20$, $65 - 15 = 50$, $74 - 24 = 50$,
$75 - 25 = 50$, $76 - 26 = 50$, $85 - 65 = 20$ e $94 - 74 = 20$ (mm). Ordenadas:
cinco valores de 20 e quatro de 50; a mediana é \qty{20}{mm.yr^{-1}}. (c) Duas
posições separadas por um ano estão na mesma fase da oscilação sazonal, que se cancela em sua diferença. Um salto afeta apenas os pares que o atravessam, aqui quatro dos nove; enquanto menos da metade dos pares atravessar um salto, a mediana o ignora; já a estimativa
pelos extremos em (a) toma o salto inteiro como movimento e sai \qty{50}{\percent} alta demais. O método falha quando os saltos são tão frequentes que a maioria dos pares de um ano contém um, ou quando a estação se move de modo irregular, por exemplo depois de um grande terremoto.
\end{solution}

\begin{solution}{pb:b1:plate-kinematics:1}
\textbf{1.} $\omega = 0.210\times(\pi/180)/\num{e6} = \qty{3.665e-9}{rad.yr^{-1}}$;
$\omega R = \num{3.665e-9}\times\qty{6.371e9}{mm} = \qty{23.35}{mm.yr^{-1}}$.

\textbf{2.} $\cos\Delta = \sin\ang{61.8}\sin\ang{64.5} + \cos\ang{61.8}
\cos\ang{64.5}\cos(\ang{-157.6}) = 0.79546 - 0.18809 = 0.60737$, logo
$\Delta = \ang{52.60}$.

\textbf{3.} $v = 23.35\sin\ang{52.60} = 23.35\times0.7944 =
\qty{18.55}{mm.yr^{-1}}$.

\textbf{4.} O numerador é $\sin\ang{157.6}\cos\ang{61.8} = 0.18011$ e o
denominador
\[
  \cos\ang{64.5}\sin\ang{61.8} - \sin\ang{64.5}\cos\ang{61.8}\cos\ang{157.6} =
  0.37940 + 0.39436 = 0.77376;
\]
ambos são positivos, logo $\beta = \arctan 0.2328 =
\ang{13.1}$. A Eurásia se move em direção a $\ang{13.1} + \ang{90} = \ang{103.1}$.

\textbf{5.} Leste: $18.55\sin\ang{103.1} = \qty{18.07}{mm.yr^{-1}}$; norte:
$18.55\cos\ang{103.1} = \qty{-4.20}{mm.yr^{-1}}$, isto é \qty{4.20}{mm.yr^{-1}}
para o sul.

\textbf{6.} $\cos\Delta = 0.8851$, $\Delta = \ang{27.74}$, $v = 23.35
\sin\ang{27.74} = \qty{10.9}{mm.yr^{-1}}$. As duas placas giram em torno de um polo na
Sibéria oriental: a dorsal de Gakkel está perto dele, e a velocidade cai com
$\sin\Delta$ embora as placas sejam as mesmas.

\textbf{7.} Oeste: $94/9.779 = \qty{9.61}{mm.yr^{-1}}$; leste: $97/9.779 =
\qty{9.92}{mm.yr^{-1}}$.

\textbf{8.} $(94 + 97)/9.779 = 191/9.779 = \qty{19.53}{mm.yr^{-1}}$.

\textbf{9.} De C2An: $54.1/2.581 = \qty{20.96}{mm.yr^{-1}}$; da base de
C1n: $19.8/0.781 = \qty{25.35}{mm.yr^{-1}}$. Elas não concordam: quanto mais jovem o
limite, maior a taxa. Uma mudança real de taxa é possível, mas um
excesso constante nas distâncias, igual em todas as idades, faria exatamente isso.

\textbf{10.} A razão entre as distâncias oeste e leste é 0.96 para 0.98: o
flanco leste cresceu alguns por cento mais. A taxa completa depende apenas da
soma das duas distâncias, logo a assimetria não a altera.

\textbf{11.} $(191 - 5.0)/9.779 = \qty{19.02}{mm.yr^{-1}}$;
$(54.1 - 5.0)/2.581 = \qty{19.02}{mm.yr^{-1}}$; $(19.8 - 5.0)/0.781 =
\qty{18.95}{mm.yr^{-1}}$. Os três agora concordam dentro de
\qty{0.1}{mm.yr^{-1}}. Os mesmos \qty{5}{km} são um quarto dos \qty{19.8}{km}
do par mais jovem, mas menos de \qty{3}{\percent} dos \qty{191}{km} do
mais antigo; limites antigos dão as taxas mais confiáveis.

\textbf{12.} A taxa de metade corrigida é $186/9.779/2 = \qty{9.51}{mm.yr^{-1}}$,
e $9.51\times0.084 = \qty{0.80}{km}$, um terço dos \qty{2.5}{km} de que
cada limite é deslocado. Seria de esperar que uma listra tão estreita fosse borrada
ao longo do caminho. Que ela seja reconhecida na dorsal de Reykjanes, enquanto está ausente na
dorsal de Gakkel, carente de magma, sugere que um suprimento abundante de
magma registra o campo fielmente mesmo em listras estreitas.

\textbf{13.} As taxas corrigidas, de 18.95 a \qty{19.02}{mm.yr^{-1}}, concordam com
os \qty{19.0}{mm.yr^{-1}} do vetor da placa; sem a correção elas o
excediam em 3 a 33 por cento.

\textbf{14.} $\cos\Delta = \sin\ang{64.139}\sin\ang{59.018} +
\cos\ang{64.139}\cos\ang{59.018}\cos\ang{27.554} = 0.97054$, logo $\Delta =
\ang{13.94}$; a corda é $2R\sin(\Delta/2) = 12742\times\sin\ang{6.97} =
\qty{1546}{km}$.

\textbf{15.} $1000/17.5 = 57$ anos.

\textbf{16.} As duas direções diferem de $\ang{123.1} - \ang{114.8} =
\ang{8.3}$: a componente horizontal ao longo do \omterm{def:b1:spherical-astronomy:spherical-triangle}{círculo máximo} é
$18.9\cos\ang{8.3} = \qty{18.70}{mm.yr^{-1}}$, e ao longo da corda
$18.70\cos\ang{6.97} = \qty{18.56}{mm.yr^{-1}}$, cerca de \qty{18.6}{mm.yr^{-1}}.

\textbf{17.} A corda se alonga \qty{17.5}{mm.yr^{-1}}, \qty{6}{\percent}
menos que os \qty{18.6}{mm.yr^{-1}} previstos. As duas estações estão longe da
fronteira de placas, logo a diferença não é local. Ou as placas se movem
um pouco mais devagar hoje do que em média nos últimos três milhões de anos,
ou uma das duas estimativas carrega um pequeno erro sistemático; as medidas
sozinhas não dizem qual. Um \omterm{def:b1:plate-kinematics:euler-pole}{vetor de Euler} calculado apenas a partir do posicionamento por satélite
prevê \qty{17.1}{mm.yr^{-1}} para esta corda, próximo do valor medido.

\textbf{18.} Leste: $21.12 - (-1.07) = \qty{22.19}{mm.yr^{-1}}$; norte:
$-6.97 - (-0.97) = \qty{-6.00}{mm.yr^{-1}}$. Velocidade $\sqrt{22.19^2 + 6.00^2} =
\qty{23.0}{mm.yr^{-1}}$, em direção a $\ang{180} - \arctan(22.19/6.00) =
\ang{105.1}$.

\textbf{19.} A direção está próxima dos \ang{103} do vetor da placa, mas
a velocidade excede seus \qty{18.6}{mm.yr^{-1}} em um quarto. As duas estações estão
dentro da zona onde a fronteira se deforma, e Höfn, perto da maior
calota, move-se com o soerguimento da \omterm{def:b1:layered-earth:layers}{crosta} à medida que o gelo afina: para cima em
\qty{12.9}{mm.yr^{-1}}, e de lado, afastando-se do gelo. As estações
dentro da zona de fronteira medem o movimento da placa mais a deformação local.

\textbf{20.} $0.210\times10 = \ang{2.10}$.

\textbf{21.} O ponto percorre um pequeno círculo de raio $R\sin\Delta$:
$6371\times0.7944\times(2.10\times\pi/180) = \qty{185.5}{km}$. A largura corrigida
da \omterm{def:b1:layered-earth:layers}{crosta} formada desde a borda jovem de C5n, \qty{186}{km} em
\qty{9.8}{Myr}, é quase a mesma: o movimento tem sido \omterm{def:b1:plate-kinematics:triple-junction}{estável} ao longo de dez
milhões de anos.

\textbf{22.} Componentes (leste, norte): Eurásia
$(22.9\sin\ang{95.1}, 22.9\cos\ang{95.1}) = (22.81, -2.04)$, Núbia
$(19.2\sin\ang{103.8}, 19.2\cos\ang{103.8}) = (18.65, -4.58)$. Eurásia em relação
à Núbia: $(4.16, 2.54)$, isto é \qty{4.9}{mm.yr^{-1}} em direção a
$\arctan(4.16/2.54) = \ang{58.6}$, lés-nordeste.

\textbf{23.} Uma junção de três dorsais é \omterm{def:b1:plate-kinematics:triple-junction}{estável} qualquer que sejam as velocidades. A
velocidade da questão 22 faz a Eurásia e a Núbia se afastarem, de lés-nordeste a
oés-sudoeste: sua fronteira na junção deve abrir-se, lentamente; é o
rifte da Terceira, cuja \omterm{def:b1:geological-maps:attitude}{direção} lés-sudeste torna a abertura oblíqua.

\textbf{24.} Parte I: \qty{18.6}{mm.yr^{-1}} no centro da Islândia
(\qty{19.0}{mm.yr^{-1}} no perfil da Parte II). Parte II: 18.95 a
\qty{19.02}{mm.yr^{-1}} depois de corrigido. Parte III: a corda Reykjavik--Stavanger
alonga-se \qty{17.5}{mm.yr^{-1}}, enquanto o vetor da placa prevê
\qty{18.6}{mm.yr^{-1}}; o par dentro da Islândia dá \qty{23.0}{mm.yr^{-1}},
o que não é o movimento da placa. O vetor da placa e as listras concordam; os
satélites diferem em \qty{6}{\percent}. As listras são médias sobre até dez
milhões de anos, o vetor da placa sobre \qty{3.16}{Myr}; os satélites medem
o movimento dos últimos vinte a trinta anos.

\textbf{25.} O vetor da placa dá \qty{18.55}{mm.yr^{-1}} através da Islândia, e
as listras corrigidas \qty{19.0}{mm.yr^{-1}} ao lado: com aproximação de um
milímetro por ano, a \omterm{def:b1:plate-kinematics:spreading-rate}{taxa de espalhamento completa} através da Islândia é
\textbf{\qty[text-series-to-math=true]{19}{mm.yr^{-1}}}.
\end{solution}
